\subsection{Numerical second-harmonic response of the dipole mode}
\label{nl:dipole-response}
We evaluated the C2 forcing in both $L=0$ and $L=2$ using separately
reconstructed numerical profiles near the certified point. The reduced
real mode components obey $\int_0^\infty(u^2+v^2)\,\dd r=1$, including
an approximate tail contribution, with $\int Y_{10}^2\,\dd\Omega=1$.
The solver uses the complex conjugate of our positive-carrier convention:
both computational unknowns have positive-$q$ outgoing Hankel conditions,
and the second unknown is conjugated to reconstruct the physical negative
sideband. The power is invariant under this convention change.

Three prescribed settings use exterior radii $48,64,64$, source spacings
$0.02,0.02,0.01$, and relative ODE tolerances
$2\times10^{-10},2\times10^{-10},2\times10^{-11}$. Two further solves at
the finest settings use the coarser reconstruction level from the same
input algorithm, with different exterior radius and tolerance. This is
an input-sensitivity comparison, not an independent implementation.
All four outgoing coefficients remain nonzero under these changes;
the largest observed relative complex-amplitude change is below
$7.8\times10^{-8}$. This sensitivity is not a rigorous error bound.

At the finest setting the harmonically averaged power coefficients are
\begin{center}
\begin{tabular}{ccc}
\toprule Angular sector&Positive sideband&Negative sideband\\\midrule
$L=0$&$2.4239\times10^{-7}$&$6.0359\times10^{-5}$\\
$L=2$&$4.4956\times10^{-7}$&$9.1587\times10^{-5}$\\\bottomrule
\end{tabular}
\end{center}
The leading power is $\epsilon^4$ times the sum of these coefficients in
the stated normalization. Both physical sidebands carry positive outward
energy with the stated outgoing conditions. These values neither determine a lifetime nor measure a
nonlinear time-evolution loss. They must not be compared directly with
coefficients using the earlier radial maximum-amplitude normalization.

This calculation supports second-harmonic leakage for the approximate C2
input, beyond merely finding open channels and nonzero local sources.
It does not certify leakage at the exact enclosed BIC. Float rounding
places some input parameters outside the tiny certified box; sampled
scaled Hermite-interpolant ODE defects have a sampled maximum of about
$5.2012\times10^{-6}$ (rounded upward to $5.3\times10^{-6}$).
These are midpoint samples on $0.01$-spaced grids extending to $r=80,96$,
excluding $|r-8|\le0.011$ around the splice, not uniform residual bounds.
Differences between reconstruction levels are much smaller. Omitted exterior
source and potential contributions were diagnosed, not converted into
propagated amplitude-error bounds in that calculation. The separate
validated monopole transfer in Section~\ref{nl:conditional-monopole}
addresses these errors using a fixed comparison function and the exact
mode's Jost normalization; it does not certify the numerical power table.

A subsequent fixed-coefficient reconstruction checked the source--current
identity for each angular sector. For the real symmetric forced operator,
$J=\operatorname{Im}(U^\dagger U')$ satisfies
$J'=\operatorname{Im}(F^\dagger U)$, where $U$ denotes the two reduced
response components and $F$ their forcing. Both computational channels
use positive-wavenumber outgoing waves. Keeping the saved matching
coefficients fixed, an independent volume quadrature agrees with
$\sum_j q_j|C_j|^2$ to relative residuals $9.6\times10^{-11}$ for $L=0$
and $1.3\times10^{-11}$ for $L=2$. This is a numerical consistency check
of the finite forced problem, not an energy-power measurement or a
validated error bound. It does not independently validate a shared source
formula. The change between the two volume quadratures is
$1.9\times10^{-9}$ or smaller relative to the respective outgoing current;
this observation is not a certified quadrature-error estimate.
